\documentclass[10pt,a4paper]{amsart}
\usepackage[margin=19mm]{geometry}
\usepackage{amsmath,amssymb,mathtools}
\usepackage[scr=rsfs,cal=euler]{mathalfa}
\usepackage{xfrac}
\usepackage[unicode,hidelinks,bookmarks=false]{hyperref}
\newcommand{\ie}{i.e.\ }
\newcommand{\cf}{cf.\ }
\newcommand{\resp}{respectively\ }
\newcommand{\Subsection}{\subsection}
\providecommand{\nobreakdash}{\nobreak\hbox{-}}
\setlength{\emergencystretch}{2em}
\multlinegap0pt
\usepackage{bm}
\usepackage{bbm}
\usepackage{dsfont}
\usepackage{xspace}
\usepackage{cancel}
\usepackage{mathtools}
\usepackage{amsthm}
\makeatletter
\@ifundefined{defn}{\newtheorem{defn}{Defn}}{}
\@ifundefined{thm}{\newtheorem{thm}[defn]{Theorem}}{}
\makeatother
\title[AVD Total Colorings of Subdivision Graphs, Joins, and Delete]{AVD Total Colorings of Subdivision Graphs, Joins, and Deleted Lexicographic Products}
\author{Amitayu Banerjee}
\date{Source: arXiv:2607.01979v3 (CC BY 4.0)}
\begin{document}
\maketitle

\noindent\emph{Source and licence.} AVD Total Colorings of Subdivision Graphs, Joins, and Deleted Lexicographic Products. Version arXiv:2607.01979v3 (CC BY 4.0), distributed under the Creative Commons Attribution 4.0 International licence, as recorded in the arXiv licence metadata for this exact version and shown on the abstract page https://arxiv.org/abs/2607.01979v3. Full rights notice: licenses/arxiv-2607.01979-CC-BY-4.0.txt.\par\smallskip

\noindent\textit{Editorial scope.} Counted unit: the Thm whose source label is \texttt{Theorem 3.6} in the article (source0.tex lines 577--579, source label \texttt{Theorem 3.6}), delivered together with the article's own complete original proof of it (source0.tex lines 581--637, one \texttt{proof} environment of 57 lines). No mathematics is written, repaired or re-derived: statement and proof are the article's own text line for line. Required context retained uncounted: none. Every definition, display and label of those lines is retained unchanged. Omitted from this package and not redistributed: 1--576, 580--580, 638--1513. Nothing inside a retained excerpt is omitted or altered: every retained body is delivered exactly as the article's lines are written, so no omission receipt of any kind accompanies this package. Citations: the retained excerpts carry no citation command at all, so no bibliography entry of the article is transcribed and no reference is redistributed. Reference adaptation: none was needed and none was made; every reference inside the delivered unit resolves inside it, so no unresolved reference or citation is produced. Printed slips kept literal: no printed word, symbol, hypothesis or number of the counted statement or of its proof was altered. Layout and class: the article's own class is \texttt{amsart} with packages including hyperref, enumerate, geometry, amsmath, amssymb, amsthm, amsfonts, xcolor, tikz, nicematrix, caption, float, verbatim, graphicx; the wrapper is the lane's pinned \texttt{amsart} editorial wrapper at 10pt on A4, which restates the theorem-like environments the retained text needs and adds the article's own macro definitions that the retained text needs (none), with extra wrapper packages (bm, bbm, dsfont, xspace, cancel, mathtools). The delivered numbering is the unit's own. Assets: no retained span contains a figure environment, a graphics-inclusion command, a PDF-page inclusion, a table or a TikZ picture; no image file is copied into this package, the package declares no asset dependency and no image file is redistributed with it. Licence: version arXiv:2607.01979v3 of this article is distributed under the Creative Commons Attribution 4.0 International licence, as recorded in the arXiv licence metadata for this exact version and shown on the abstract page https://arxiv.org/abs/2607.01979v3. A full rights notice is delivered at licenses/arxiv-2607.01979-CC-BY-4.0.txt. Characters: every retained excerpt is pure ASCII, so the delivered engine, XeLaTeX with its default Latin Modern text font, reports no missing character.
\begin{thm}\label{Theorem 3.6}
If $G$ is a graph of order $m$ such that $m\geq \Delta(G)+3$ and $\chi''_{a}(G)\leq \Delta(G)+2$, then $\chi''_{a}(D_{lex}(P_n, G))\leq \Delta(D_{lex}(P_n, G))+2$ for every  $n\geq 3$. 
\end{thm}

\begin{proof}
Fix an integer $n\geq 3$. Let $V(P_n)=\{v_1,...,v_n\}$ be an enumeration of the set of vertices of $P_n$. In $D_{lex}(P_n, G)$, there are $n$ copies of $G$ and the maximum degree of $D_{lex}(P_n, G)$ is $\Delta(G)+2(m-1)$. 
Let $G_{v_i}$ denote the copy of $G$ with respect to the vertex $v_i\in V(P_n)$.
Since $m-1\geq\Delta(G)+2$, we split $\Delta(G)+2(m-1)+2$ colors into the following disjoint sets:
$S_1=\{a_1,...,a_{\Delta(G)+2}\}$, $S_2=\{b_1,...,b_{\Delta(G)+2}\}$, $S'_{2}=\{b_{\Delta(G)+3},...,b_{m-1}\}$, $S_3=\{c_1,...,c_{\Delta(G)+2}\}$, $S'_{3}=\{c_{\Delta(G)+3},...,c_{m-1}\}$.\footnote{If $m-1=\Delta(G)+2$ then $S'_2=\emptyset$ and $S'_3=\emptyset$.} Define a total coloring
$c:V(D_{lex}(P_n,G))\cup E(D_{lex}(P_n,G))
\longrightarrow \bigcup_{i=1}^{3}S_i\cup S'_2\cup S'_3$ as follows: 

\begin{enumerate}
    \item Since $\chi''_{a}(G)\leq \Delta(G)+2$, consider an AVD-total coloring of $G_{v_1}$ with colors in $S_1$. 
    The edges between $G_{v_{1}}$ and $G_{v_{2}}$ are colored with the colors in $S_2\cup S'_2$.
    \item Similar to (1), consider 
    an AVD-total coloring of $G_{v_2}$ with colors in $S_3$ and color the edges between
    $G_{v_{2}}$ and $G_{v_{3}}$ using the colors in $S'_3\cup S_1$. 
    
    \item Assign an AVD-total coloring to $G_{v_3}$ using colors from the set $S_2$
    and assign the colors in the set $S'_2\cup S_3$ to the edges between $G_{v_3}$ and $G_{v_4}$.
    
    \item Consider 
    an AVD-total coloring of $G_{v_4}$ with colors in $S_1$ and color the edges between $G_{v_4}$ and $G_{v_5}$ with the colors in the set $S_2\cup S'_3$.

    \item Using the colors in the set $S_3$, consider 
    an AVD-total coloring of $G_{v_5}$
    and use colors in the set $S_1\cup S'_2$ for the edges between $G_{v_5}$ and $G_{v_6}$.

    \item Assign an AVD-total coloring to $G_{v_6}$ using colors from the set $S_2$ and assign the colors in the set $S_3\cup S'_3$ to the edges between $G_{v_6}$ and $G_{v_7}$.

    \item Color the vertices and edges of $G_{v_{7}}$ as in $G_{v_{1}}$, and repeat the previous steps till all the elements of $D_{lex}(P_n, G)$ are colored.
\end{enumerate}

It is clear that the coloring $c$ is a proper total coloring of $D_{lex}(P_n,G)$. 
We show that $c$ is an AVD-total coloring of $D_{lex}(P_n,G)$. 
Let $xy\in E(D_{lex}(P_n,G))$. 

\textbf{Case (i).} Suppose $x=(v_i,a)\in V(G_{v_i})$ and $y=(v_i,b)\in V(G_{v_i})$ for some $a,b\in V(G)$ and $v_{i} \in V(P_n)$.
Since we considered an AVD-total coloring of $G_{v_i}$, we obtain $C_{G_{v_i}}(x)\neq C_{G_{v_i}}(y)$.  Moreover, $C_{G_{v_i}}(x), C_{G_{v_i}}(y)\subseteq S_k$ for some $k\in \{1,2,3\}$.
Without loss of generality, let $k=1$. Then
\[
C_{D_{lex}(P_n, G)}(x)\setminus C_{G_{v_i}}(x),\;
C_{D_{lex}(P_n, G)}(y)\setminus C_{G_{v_i}}(y)
\subseteq S_2\cup S'_2\cup S_3 \cup S'_3,
\]
where $S_1 \cap (S_2\cup S'_2\cup S_3 \cup S'_3)=\emptyset$.
Consequently, $C_{D_{lex}(P_n, G)}(x)\neq C_{D_{lex}(P_n, G)}(y)$.
 
\textbf{Case (ii).} Suppose $x=(v_i,a)\in V(G_{v_i})$ and $y=(v_j,b)\in V(G_{v_j})$ for some $a,b\in V(G)$ and $v_i,v_j\in V(P_n)$ such that $i< j<n$. By the definition of $c$, the elements of $G_{v_i}$ are colored with the colors in $S_k$ for some $k\in\{1,2,3\}$.
Among the incident edges of $x$ in $D_{lex}(P_n, G)$, only the edges of $G_{v_i}$ incident to $x$ are colored with the colors in $S_k$. Since $x$ is also colored with the colors in $S_k$, we obtain $|S_k\cap C_{D_{lex}(P_n, G)}(x)|\leq \Delta(G)+1$.
However, $v_iv_j\in E(P_n)$ since $xy\in E(D_{lex}(P_n,G))$. Thus, $j=i+1$ as we assumed $i<j$. 
By the definition of $c$, all the colors in $S_k$ are used to color the edges between $G_{v_j}$ and $G_{v_{j+1}}$.
Consequently, 
\[|S_k\cap C_{D_{lex}(P_n, G)}(y)|= |S_k| = \Delta(G)+2>|S_k\cap C_{D_{lex}(P_n, G)}(x)|.\] 
Thus, there exists $a\in S_k$ such that $a\in C_{D_{lex}(P_n, G)}(y)$ but  $a\not\in C_{D_{lex}(P_n, G)}(x)$. Hence $C_{D_{lex}(P_n, G)}(x)\neq C_{D_{lex}(P_n, G)}(y)$.


\textbf{Case (iii).} Suppose $x=(v_{n-1},a)\in V(G_{v_{n-1}})$ and $y=(v_n,b)\in V(G_{v_n})$ for some $a,b\in V(G)$. If the elements of $G_{v_n}$ are colored with the colors in $S_k$ for some $k\in\{1,2,3\}$, then similar to the arguments of Case (ii), we have $|S_k\cap C_{D_{lex}(P_n, G)}(y)|\leq \Delta(G)+1$ whereas $|S_k\cap C_{D_{lex}(P_n, G)}(x)|= \Delta(G)+2$. Thus, $C_{D_{lex}(P_n, G)}(x)\neq C_{D_{lex}(P_n, G)}(y)$.
%\textbf{Case (II):} Suppose $m=\Delta(G)+3$. Assume $S_1,S_2,S_2',S_3,S'_3$ as in Case (I). We note that $S'_2=\emptyset$ and $S'_3=\emptyset$ in this case. The rest of the proof proceeds as in Case (I).
\end{proof}

\end{document}
